\documentclass[12pt,a4paper]{article}
\usepackage[utf8]{inputenc}
\usepackage{amsmath}
\usepackage{amsfonts}
\usepackage{amssymb}
\usepackage{lmodern}
\usepackage{amsmath}
\usepackage{amsthm}
\usepackage{enumerate}
\usepackage{graphicx}
\usepackage{mathtools}
\usepackage{stackrel}
\usepackage[left=2cm,right=2cm,top=2cm,bottom=2cm]{geometry}
\newcounter{neq}
\providecommand{\abs}[1]{\lvert#1\rvert}
\newcommand{\SIGMA}{\Sigma^{\ast}}
\newcommand{\PN}{\par\noindent}

\newtheorem{theorem}[equation]{Theorem}
\newtheorem{lemma}[equation]{Lemma}
\newtheorem{proposition}[equation]{Proposition}
\newtheorem{corollary}[equation]{Corollary}

\author{Agustín Curto, agucurto95@gmail.com \\
			 Francisco Nievas, frannievas@gmail.com}
\title{Resumen de teoremas para el final \\ de Lógica}
\date{2017}

\begin{document}
	\clearpage\maketitle
 	\thispagestyle{empty}
 	\tableofcontents

	\vspace{5cm}
	\PN \textbf{Nota:} Este resumen se corresponde con la materia dictada en el año 2017. Los autores no se
	responsabilizan de posibles cambios que pudiesen realizarse en los temas dictados en la misma, así como tampoco de
	errores involuntarios que pudiesen existir en dicho resumen.

	\vspace{\fill}
	\begin{center}
		Por favor, mejorá este documento en github
		\includegraphics[width=1cm]{graphics/github.png} \\
		https://github.com/ResumenesFaMAF/resumenLogica
	\end{center}

\end{document}
